\begin{abcliste}
\abc
Bei konstanter Geschwindigkeit ist die Kraft so gross wie die Gewichtskraft:
\begin{align}
 F_\mathrm{A} &= \FAF = \mK\cdot\g = \FA \approx \result{\FAP{0}{3}}\\
 W_\mathrm{A} &= \WAF = \mK\cdot\g\cdot\hK = \WA \approx \result{\WAP{0}{3}}
\end{align}

\abc
Auf der Rampe muss man nur die Komponente der Gewichtskraft entlang der Rampe aufbringen, $m\,g\sin\alpha$ mit $\sin\alpha = h/L$:
\begin{align}
 F_\mathrm{B} &= \FBF = \frac{\mK\cdot\g\cdot\hK}{\LR} = \FB \approx \result{\FBP{0}{3}}\\
 W_\mathrm{B} &= \WBF = \FB\cdot\LR = \WB \approx \result{\WBP{0}{3}}
\end{align}

\abc
Die Rampe ist $L/h = \nR \approx \nRP{0}{3}$-mal so lang wie hoch. Um genau diesen Faktor ist die nötige Kraft kleiner als $m\,g$, der Weg aber um denselben Faktor länger. Das Produkt Kraft mal Weg bleibt gleich:
\begin{align}
 W_\mathrm{B} &= \frac{m\,g\,h}{L}\cdot L = m\,g\,h = W_\mathrm{A}
\end{align}
Die Arbeit erhöht in beiden Fällen nur die Lageenergie der Kiste um $m\,g\,h$, und die hängt allein von der Höhe ab, nicht vom Weg. Die Rampe spart Kraft, aber keine Arbeit.
\end{abcliste}
